\documentclass[11pt,a4paper]{article}
% =========================================================
%                   PAQUETES PRINCIPALES
% =========================================================
\usepackage[spanish, es-tabla]{babel}
\usepackage[utf8]{inputenc}
\usepackage[T1]{fontenc}
\usepackage{graphicx}
\usepackage{float}
\usepackage{geometry}
\usepackage{setspace}
\usepackage{hyperref}
\usepackage{tocloft}
\usepackage{caption}
\usepackage{array}
\usepackage{booktabs}
\usepackage{ragged2e}
\usepackage{enumitem}
\usepackage{listings}
\usepackage{xcolor}
\usepackage{amsmath}
\usepackage{amssymb}
\usepackage{tikz}
\usepackage{pgfplots}
\usepackage{multirow}
\usepackage{longtable}
\usepackage{fancyhdr}
\usepackage{mdframed}
\usepackage{subcaption}
\usepackage{wrapfig}
\usepackage{url}
\usepackage{titlesec}
\usepackage{lastpage}
\usepackage{tabularx}
\usetikzlibrary{positioning, arrows.meta, shapes.geometric, calc}
\pgfplotsset{compat=1.18}

% =========================================================
%                  CONFIGURACIÓN GENERAL
% =========================================================
\geometry{
    top=2.8cm,
    bottom=2.8cm,
    left=3cm,
    right=3cm
}
\setstretch{1.35}
\captionsetup{
    font=small,
    labelfont=bf
}
\hypersetup{
    colorlinks=true,
    linkcolor=black,
    citecolor=blue,
    urlcolor=blue,
    pdftitle={GIMNASIO: GYMsos - Operating System Inteligente para Gimnasios},
    pdfauthor={Paipay Vega, Eduardo Sebastian}
}
\shorthandoff{"}

% =========================================================
%                        COLORES
% =========================================================
\definecolor{bluekeywords}{rgb}{0.13,0.13,1}
\definecolor{greencomments}{rgb}{0,0.5,0}
\definecolor{redstrings}{rgb}{0.9,0,0}
\definecolor{graynumbers}{rgb}{0.5,0.5,0.5}
\definecolor{lightgray}{rgb}{0.96,0.96,0.96}
\definecolor{darkblue}{RGB}{0,51,102}
\definecolor{accent}{RGB}{0,102,204}
\definecolor{success}{RGB}{46,125,50}
\definecolor{warning}{RGB}{251,140,0}

% =========================================================
%                  CONFIGURACIÓN DE CÓDIGO
% =========================================================
\lstset{
    language=bash,
    showspaces=false,
    showtabs=false,
    breaklines=true,
    showstringspaces=false,
    breakatwhitespace=true,
    commentstyle=\color{greencomments},
    keywordstyle=\color{bluekeywords}\bfseries,
    stringstyle=\color{redstrings},
    basicstyle=\ttfamily\footnotesize,
    frame=single,
    backgroundcolor=\color{lightgray},
    numbers=left,
    numberstyle=\tiny\color{graynumbers},
    stepnumber=1,
    tabsize=4,
    captionpos=b
}

% =========================================================
%                 ESTILO DE TÍTULOS
% =========================================================
\titleformat{\section}
{\normalfont\Large\bfseries\color{darkblue}}
{\thesection.}{0.5em}{}
\titleformat{\subsection}
{\normalfont\large\bfseries\color{accent}}
{\thesubsection.}{0.5em}{}
\titleformat{\subsubsection}
{\normalfont\normalsize\bfseries}
{\thesubsubsection.}{0.5em}{}

% =========================================================
%                ENCABEZADO Y PIE DE PÁGINA
% =========================================================
\pagestyle{fancy}
\fancyhf{}
\fancyhead[L]{\small\textit{GYMsos - Operating System para Gimnasios}}
\fancyhead[R]{\small\textit{Proyecto GIMNASIO}}
\fancyfoot[L]{\small Documento Técnico v3.0}
\fancyfoot[R]{\small Página \thepage\ de \pageref{LastPage}}
\renewcommand{\headrulewidth}{0.4pt}
\renewcommand{\footrulewidth}{0.4pt}

% =========================================================
%               ENTORNO DE DEFINICIONES
% =========================================================
\newmdenv[
  backgroundcolor=blue!5,
  linecolor=accent,
  linewidth=1.5pt,
  leftmargin=0pt,
  rightmargin=0pt,
  innertopmargin=8pt,
  innerbottommargin=8pt,
  innerleftmargin=10pt,
  innerrightmargin=10pt
]{definicion}

\newmdenv[
  backgroundcolor=green!5,
  linecolor=success,
  linewidth=1.5pt,
  leftmargin=0pt,
  rightmargin=0pt,
  innertopmargin=8pt,
  innerbottommargin=8pt,
  innerleftmargin=10pt,
  innerrightmargin=10pt
]{beneficio}

\newmdenv[
  backgroundcolor=orange!5,
  linecolor=warning,
  linewidth=1.5pt,
  leftmargin=0pt,
  rightmargin=0pt,
  innertopmargin=8pt,
  innerbottommargin=8pt,
  innerleftmargin=10pt,
  innerrightmargin=10pt
]{alerta}

% =========================================================
%                       DOCUMENTO
% =========================================================
\begin{document}

% =========================================================
%                         PORTADA
% =========================================================
\begin{titlepage}
\centering
\vspace*{1cm}
{\Large \textbf{PROYECTO GIMNASIO}}\\[0.4cm]
{\large Documento Técnico y Estratégico}\\[1cm]
\vspace{1cm}
{\Huge \textbf{GYMsos}}\\[0.5cm]
\rule{15cm}{0.8pt}\\[0.4cm]
{\LARGE \textbf{Operating System Inteligente para Gimnasios}}\\[0.4cm]
\rule{15cm}{0.8pt}
\vspace{2cm}
{\large \textbf{Documento Técnico - Versión 3.0}}\\[1.5cm]
{\large \textbf{Especificaciones Funcionales, Requisitos y Arquitectura}}\\[2cm]
\begin{table}[H]
\centering
\renewcommand{\arraystretch}{1.5}
\begin{tabular}{|p{5cm}|p{7cm}|}
\hline
\textbf{Estudiante} & Paipay Vega, Eduardo Sebastian \\ \hline
\textbf{Código} & 27222136 \\ \hline
\textbf{Universidad} & Universidad Nacional de San Cristóbal de Huamanga \\ \hline
\textbf{Versión} & 3.0 - UNICORN LEVEL \\ \hline
\textbf{Fecha} & Mayo 2026 \\ \hline
\textbf{Estado} & 100\% Completado \\ \hline
\end{tabular}
\end{table}
\vfill
{\large Ayacucho -- Perú}\\[0.2cm]
{\large 2026}
\end{titlepage}

% =========================================================
%              TABLA DE CONTENIDOS
% =========================================================
\newpage
\tableofcontents
\newpage

% =========================================================
%                RESUMEN EJECUTIVO
% =========================================================
\section*{Resumen Ejecutivo}
\addcontentsline{toc}{section}{Resumen Ejecutivo}

\subsection*{Visión General}

GYMsos es el \textbf{Operating System inteligente para gimnasios} — una plataforma integral que transforma la gestión tradicional basada en Excel hacia un ecosistema digital automatizado, inteligente y escalable.

\subsection*{Objetivos Principales}

\begin{itemize}
    \item Automatizar procesos administrativos de gestión de gimnasios
    \item Proporcionar análisis de datos en tiempo real
    \item Mejorar experiencia de usuario mediante interfaz moderna
    \item Integrar múltiples canales de comunicación
    \item Permitir escalabilidad a múltiples sucursales
    \item Implementar control de acceso moderno (QR + biometría)
\end{itemize}

\subsection*{La Solución: 13 Innovaciones Tecnológicas}

GYMsos integra 13 innovaciones tecnológicas:

\begin{enumerate}
    \item \textbf{Netflix Fitness} — Motor de recomendaciones con IA basado en historial
    \item \textbf{Churn AI} — Sistema de detección predictiva de patrones
    \item \textbf{Digital Twin} — Representación virtual del gimnasio en tiempo real
    \item \textbf{Gamificación Masiva} — Sistema de puntuación y competencias
    \item \textbf{Social Network} — Red social integrada de miembros
    \item \textbf{Smart Gym OS} — Control inteligente de espacios y máquinas
    \item \textbf{Autonomous Growth Engine} — Sistema de crecimiento automático
    \item \textbf{Preventive Health} — Análisis de salud con integración de wearables
    \item \textbf{B2B2C Corporate} — Módulo de gestión empresarial
    \item \textbf{AI Copilot para Trainers} — Asistente inteligente para entrenadores
    \item \textbf{Spotify Recommendations} — Integración de música personalizada
    \item \textbf{Marketplace} — Plataforma integrada de productos y servicios
    \item \textbf{Tesla Moment} — Experiencia de usuario premium
\end{enumerate}

% =========================================================
%              DESCRIPCIÓN DEL PROYECTO
% =========================================================
\section{Descripción del Proyecto}

\subsection{¿Qué es GYMsos?}

GYMsos no es solo un software de gestión de gimnasios. Es un \textbf{ecosistema completo} que integra:

\begin{itemize}
    \item \textbf{Aplicación Web} — Dashboard para gerentes y administradores
    \item \textbf{Aplicación Móvil} — App para miembros (iOS, Android)
    \item \textbf{Aplicación Desktop} — Herramientas para recepcionistas y staff
    \item \textbf{Integración WhatsApp} — Comunicación automática con miembros
    \item \textbf{Control de Acceso} — QR + biometría integrada
    \item \textbf{Análisis de Datos} — Dashboards en tiempo real con IA
    \item \textbf{Marketplace} — Plataforma de productos y servicios complementarios
    \item \textbf{API Abierta} — Para integraciones con terceros
\end{itemize}

\subsection{Actores del Sistema}

\begin{table}[H]
\centering
\caption{Actores Principales del Sistema}
\begin{tabular}{|l|p{6cm}|l|}
\hline
\textbf{Actor} & \textbf{Funcionalidades Acceso} & \textbf{Plataforma} \\ \hline
Miembro & Ver perfil, historial, logros, marketplace & Mobile App \\ \hline
Recepcionista & Gestión de acceso, atención al cliente & Desktop + Web \\ \hline
Entrenador & Visualizar clientes, recomendaciones & Mobile + Dashboard \\ \hline
Gerente & Dashboards, reportes, configuración & Web Dashboard \\ \hline
Administrador & Control total del sistema & Web Admin Panel \\ \hline
\end{tabular}
\end{table}

\subsection{Especificaciones de Escala}

\begin{itemize}
    \item \textbf{Capacidad}: 5,000+ miembros por sucursal
    \item \textbf{Arquitectura}: Microservicios distribuidos
    \item \textbf{Base de datos}: PostgreSQL con replicación
    \item \textbf{Caché}: Redis para sesiones y datos frecuentes
    \item \textbf{Disponibilidad}: 99.9\% uptime SLA
    \item \textbf{Latencia}: < 2 segundos en 95\% de requests
    \item \textbf{Backups}: Diarios, con RTO < 4 horas
\end{itemize}

% =========================================================
%            ANÁLISIS DE REQUISITOS INICIALES
% =========================================================
\section{Análisis de Requisitos Iniciales}

\subsection{Necesidades del Sistema}

\subsubsection{Gestión Automatizada}

El sistema debe automatizar:
\begin{itemize}
    \item Registro de membresías
    \item Renovaciones de membresías
    \item Control de pagos
    \item Reportes de asistencia
    \item Generación de reportes administrativos
\end{itemize}

\subsubsection{Control de Acceso}

Implementar:
\begin{itemize}
    \item Control de acceso mediante QR
    \item Integración con biometría
    \item Trazabilidad completa de accesos
    \item Integración con datos de miembros
    \item Análisis de patrones de uso
\end{itemize}

\subsubsection{Análisis de Datos}

El sistema debe proporcionar:
\begin{itemize}
    \item Dashboards en tiempo real
    \item Análisis predictivo de patrones
    \item Reportes automáticos
    \item Segmentación de usuarios
    \item Estadísticas de uso
\end{itemize}

\subsubsection{Escalabilidad Multi-Sucursal}

Arquitectura que soporte:
\begin{itemize}
    \item Múltiples sucursales con gestión centralizada
    \item Sincronización de datos entre ubicaciones
    \item Gestión de permisos y roles por sucursal
    \item Reportes consolidados
\end{itemize>

\subsection{Mapa de Soluciones}

\begin{table}[H]
\centering
\caption{Requisitos vs Módulos de Solución}
\begin{tabularx}{\textwidth}{|l|X|l|}
\hline
\textbf{\#} & \textbf{Requisito} & \textbf{Módulo} \\ \hline
1 & Gestión automatizada & Member Service \\ \hline
2 & Control de acceso & Attendance Service \\ \hline
3 & Análisis de datos & Analytics Service \\ \hline
4 & Multi-sucursal & Architecture Centralized \\ \hline
5 & Notificaciones & Notification Service \\ \hline
6 & Integración pagos & Payment Service \\ \hline
7 & Marketplace & Marketplace Service \\ \hline
8 & IA y Predicción & AI Service \\ \hline
\end{tabularx}
\end{table}

% =========================================================
%              LAS 13 INNOVACIONES
% =========================================================
\section{Las 13 Disruptivas Innovaciones}

\subsection{1. Netflix Fitness — Motor de Recomendaciones}

\begin{definicion}
\textbf{Concepto}: Sistema de recomendación de ejercicios basado en IA que analiza historial y preferencias.
\end{definicion}

\textbf{Componentes técnicos}:
\begin{itemize}
    \item Base de datos de preferencias y historial de usuario
    \item Algoritmo de recomendación basado en machine learning
    \item Tracking de progresión y feedback
    \item API de recomendaciones en tiempo real
\end{itemize}

\textbf{Integraciones}: AI Service, Member Service

\subsection{2. Churn AI — Detección Predictiva}

\begin{definicion}
\textbf{Concepto}: Sistema de predicción basado en análisis de patrones de comportamiento
\end{definicion}

\textbf{Técnicas implementadas}:
\begin{itemize}
    \item Análisis de patrones de asistencia
    \item Machine learning para detección de anomalías
    \item Alertas automáticas basadas en modelos predictivos
    \item Sistema de notificaciones integrado
\end{itemize}

\textbf{Integraciones}: Analytics Service, Notification Service

\subsection{3. Digital Twin — Representación Virtual}

\begin{definicion}
\textbf{Concepto}: Modelo virtual sincronizado con estado del gimnasio
\end{definicion}

\textbf{Componentes}:
\begin{itemize}
    \item Estado en tiempo real de máquinas y espacios
    \item Sensor de ocupación por zona
    \item Disponibilidad de equipamiento
    \item Base de datos de inventario
\end{itemize}

\subsection{4. Gamificación — Sistema de Puntuación}

\begin{definicion}
\textbf{Concepto}: Sistema de gamificación integrado en la plataforma
\end{definicion}

\textbf{Elementos técnicos}:
\begin{itemize}
    \item Módulo de cálculo de puntos y logros
    \item Sistema de leaderboards en base de datos
    \item Gestión de badges y recompensas
    \item APIs para interacción gamificada
\end{itemize}

\subsection{5. Social Network — Red Social Integrada}

\begin{definicion}
\textbf{Concepto}: Funcionalidades de red social dentro de la aplicación
\end{definicion}

\textbf{Funcionalidades}:
\begin{itemize}
    \item Sistema de seguimiento de usuarios
    \item Feed de actividades
    \item Mensajería privada
    \item Grupos y desafíos colaborativos
\end{itemize}

\subsection{6. Smart Gym OS — Control de Espacios}

\begin{definicion}
\textbf{Concepto}: Integración de IoT y control automatizado de recursos
\end{definicion}

\textbf{Integraciones IoT}:
\begin{itemize}
    \item APIs para máquinas conectadas
    \item Sistemas de control de ambiente (HVAC, iluminación)
    \item Sensores de ocupación
    \item Sistema de reservas en tiempo real
\end{itemize}

\subsection{7. Autonomous Growth Engine — Sistema de Referrals}

\begin{definicion}
\textbf{Concepto}: Sistema automatizado de referrals y crecimiento
\end{definicion}

\textbf{Implementación}:
\begin{itemize}
    \item Módulo de invitaciones automáticas
    \item Tracking de referrals en BD
    \item Sistema de incentivos automático
    \item Análisis de viral loops
\end{itemize}

\subsection{8. Preventive Health — Integración Biométrica}

\begin{definicion}
\textbf{Concepto}: Integración con wearables y dispositivos de salud
\end{definicion}

\textbf{Integraciones}:
\begin{itemize}
    \item APIs de Apple Health, Google Fit
    \item Sincronización de datos biométricos
    \item Análisis de métricas de salud
    \item Alertas basadas en umbrales
\end{itemize}

\subsection{9. B2B2C Corporate — Módulo Empresarial}

\begin{definicion}
\textbf{Concepto}: Funcionalidades para gestión de programas corporativos
\end{definicion}

\textbf{Características técnicas}:
\begin{itemize}
    \item Dashboard dedicado para administradores corporativos
    \item Gestión de empleados en lotes
    \item Reportes específicos para empresas
    \item Sistema de autorización multi-nivel
\end{itemize}

\subsection{10. AI Copilot para Trainers — Asistente IA}

\begin{definicion}
\textbf{Concepto}: Sistema de recomendación asistido por IA para entrenadores
\end{definicion}

\textbf{Tecnologías}:
\begin{itemize}
    \item Modelos de IA para personalización
    \item Computer vision para análisis de forma
    \item Generación automática de reportes
    \item Integración con historial de entrenamientos
\end{itemize}

\subsection{11. Spotify Recommendations — Integración Musical}

\begin{definicion}
\textbf{Concepto}: Integración con servicios de música para crear listas personalizadas
\end{definicion}

\textbf{Implementación}:
\begin{itemize}
    \item APIs de Spotify/Apple Music
    \item Algoritmo de matching música-ejercicio
    \item Historial de preferencias musicales
    \item Generación de playlists basadas en tempo
\end{itemize}

\subsection{12. Marketplace — Plataforma de E-commerce}

\begin{definicion}
\textbf{Concepto}: Plataforma integrada de productos y servicios
\end{definicion}

\textbf{Componentes técnicos}:
\begin{itemize}
    \item Catálogo de productos en BD
    \item Carrito de compras
    \item Integración con Payment Service
    \item Sistema de pedidos y tracking
\end{itemize}

\subsection{13. Tesla Moment — Experiencia de Usuario}

\begin{definicion}
\textbf{Concepto}: Interfaz moderna y fluida con experiencia premium
\end{definicion}

\textbf{Aspectos técnicos}:
\begin{itemize}
    \item Diseño responsive (web + mobile)
    \item Animaciones y transiciones fluidas
    \item Personalización de interfaz por usuario
    \item Integración con sensores biométricos
\end{itemize}

% =========================================================
%              ESPECIFICACIONES FUNCIONALES
% =========================================================
\section{Especificaciones Funcionales}

\subsection{Requisitos Funcionales Principales (RF)}

GYMsos debe implementar los siguientes requisitos funcionales:

\subsubsection{Módulo de Gestión de Membresías}

\begin{itemize}
    \item RF-001: Registrar nuevo miembro con datos personales
    \item RF-002: Actualizar información de miembro
    \item RF-003: Renovar membresía automáticamente
    \item RF-004: Registrar tipo de membresía (mensual, trimestral, anual)
    \item RF-005: Historial de cambios de membresía
\end{itemize}

\subsubsection{Módulo de Control de Acceso}

\begin{itemize}
    \item RF-010: Generar código QR único por miembro
    \item RF-011: Validar QR en entrada del gimnasio
    \item RF-012: Registrar hora y fecha de asistencia
    \item RF-013: Integración con sistemas biométricos
    \item RF-014: Auditoría completa de accesos
\end{itemize}

\subsubsection{Módulo de Analytics}

\begin{itemize}
    \item RF-020: Generar dashboard de ocupación en tiempo real
    \item RF-021: Reporte de asistencia por período
    \item RF-022: Análisis de máquinas más utilizadas
    \item RF-023: Horas pico de ocupación
    \item RF-024: Exportar reportes en PDF/Excel
\end{itemize}

\subsubsection{Módulo de Notificaciones}

\begin{itemize}
    \item RF-030: Enviar notificaciones push a miembros
    \item RF-031: Enviar emails automáticos
    \item RF-032: Integración con WhatsApp
    \item RF-033: Gestión de preferencias de notificación
    \item RF-034: Log de notificaciones enviadas
\end{itemize}

\subsubsection{Módulo de Gamificación}

\begin{itemize}
    \item RF-040: Calcular puntos por asistencia
    \item RF-041: Mantener leaderboards
    \item RF-042: Gestionar badges y logros
    \item RF-043: Sistema de competencias
    \item RF-044: Historial de puntos del usuario
\end{itemize}

\subsubsection{Módulo de Marketplace}

\begin{itemize}
    \item RF-050: Catálogo de productos
    \item RF-051: Carrito de compras
    \item RF-052: Procesamiento de transacciones
    \item RF-053: Historial de compras
    \item RF-054: Disponibilidad de inventario
\end{itemize}

\subsection{Requisitos No Funcionales (RNF)}

\begin{itemize}
    \item RNF-001: Disponibilidad: 99.9\% uptime
    \item RNF-002: Rendimiento: Carga de página < 2 segundos
    \item RNF-003: Escalabilidad: Soportar 5,000+ miembros
    \item RNF-004: Seguridad: Encriptación end-to-end
    \item RNF-005: Cumplimiento: GDPR, normativas locales
    \item RNF-006: Compatibilidad: Web, iOS, Android
    \item RNF-007: Respaldo: Backup automático diario
    \item RNF-008: Recuperación: RTO < 4 horas
\end{itemize}

% =========================================================
%            ARQUITECTURA TÉCNICA
% =========================================================
\section{Arquitectura Técnica}

\subsection{Stack Tecnológico}

\begin{table}[H]
\centering
\caption{Tecnologías Utilizadas}
\begin{tabular}{|l|l|}
\hline
\textbf{Componente} & \textbf{Tecnología} \\ \hline
Frontend Web & React 18, TypeScript, Tailwind CSS \\ \hline
Frontend Mobile & React Native / Flutter \\ \hline
Backend API & Node.js + Express / Python + FastAPI \\ \hline
Base de Datos & PostgreSQL (relacional), Redis (caché) \\ \hline
IA/ML & TensorFlow, scikit-learn, OpenAI API \\ \hline
Autenticación & Auth0, OAuth 2.0, JWT \\ \hline
Pagos & Stripe, PayU, MercadoPago APIs \\ \hline
Comunicaciones & Twilio (WhatsApp), Firebase (push) \\ \hline
Hosting & AWS (EC2, S3, RDS), Docker, Kubernetes \\ \hline
\end{tabular}
\end{table}

\subsection{Arquitectura de Microservicios}

La plataforma está dividida en servicios independientes:

\begin{itemize}
    \item \textbf{Authentication Service} — Gestión de usuarios y permisos
    \item \textbf{Member Service} — CRUD de miembros y membresías
    \item \textbf{Payment Service} — Procesamiento de pagos integrado
    \item \textbf{Attendance Service} — Control de acceso y asistencia
    \item \textbf{Analytics Service} — Dashboards y reportes IA
    \item \textbf{Notification Service} — Email, SMS, push, WhatsApp
    \item \textbf{Marketplace Service} — Plataforma de e-commerce
    \item \textbf{AI Service} — Modelos predictivos, recomendaciones
\end{itemize}

\subsection{Seguridad}

\begin{alerta}
\textbf{Cumplimiento}: GDPR, CCPA, normativas locales de datos
\end{alerta}

\begin{itemize}
    \item Encriptación end-to-end para datos sensibles
    \item Auditoría de accesos completamente logueada
    \item Backup automático y disaster recovery
    \item Penetration testing trimestral
    \item Certificados SSL/TLS para todas las comunicaciones
\end{itemize}

% =========================================================
%              ROADMAP Y FASES
% =========================================================
\section{Roadmap de Implementación}

\subsection{Fase 1: Problemas y Oportunidades (Completada)}

\textbf{Objetivo}: Identificar y cuantificar problemas del mercado

\begin{itemize}
    \item Análisis de gestión manual e ineficiencia
    \item Cuantificación de churn y costo
    \item Identificación de stakeholders
    \item Definición de restricciones
\end{itemize}

\subsection{Fase 2: Valor Agregado (Completada)}

\textbf{Objetivo}: Definir 13 innovaciones que crean valor

\begin{itemize}
    \item Netflix Fitness, Churn AI, Digital Twin
    \item Gamificación, Social Network
    \item Smart Gym OS, Growth Engine
    \item 6 innovaciones más
\end{itemize}

\subsection{Fase 3: Requisitos Funcionales y Casos de Uso (Próxima)}

\textbf{Objetivo}: Especificar exactamente qué construir

\begin{itemize}
    \item 50+ Requisitos Funcionales
    \item 30+ Casos de Uso
    \item Diagramas de flujo
    \item Especificación de APIs
\end{itemize}

\subsection{Fase 4: Plan de Negocio (Completada)}

\textbf{Objetivo}: Modelo de negocio y proyecciones

\begin{itemize}
    \item Modelo de ingresos multi-stream
    \item Proyecciones 3 años
    \item Análisis de TAM/SAM/SOM
    \item ROI y breakeven
\end{itemize}

\subsection{Fase 5: Diseño de Base de Datos (Próxima)}

\textbf{Objetivo}: Estructura de datos optimizada

\begin{itemize}
    \item Tablas normalizadas
    \item Índices y relaciones
    \item Esquema para 5000+ miembros
\end{itemize}

\subsection{Fase 6: UX/UI Design (Próxima)}

\textbf{Objetivo}: Experiencia de usuario intuitiva

\begin{itemize}
    \item Wireframes de 25+ pantallas
    \item Flujos de usuario
    \item Guía de diseño
    \item Prototipo interactivo
\end{itemize}

\subsection{Fase 7: Desarrollo e Implementación (Próxima)}

\textbf{Objetivo}: Construcción del producto

\begin{itemize}
    \item MVP con 5 innovaciones core
    \item Beta testing con 3-5 gimnasios
    \item Iteración rápida
    \item Lanzamiento público
\end{itemize}

% =========================================================
%              CASOS DE USO (CU)
% =========================================================
\section{Casos de Uso Principales}

\subsection{Actor: Miembro}

\subsubsection{CU-001: Registrarse en la Plataforma}

\begin{definicion}
\textbf{Descripción}: El miembro crea su cuenta en GYMsos proporcionando datos personales.

\textbf{Flujo Principal}:
\begin{itemize}
    \item Actor accede a la aplicación
    \item Sistema presenta formulario de registro
    \item Actor ingresa email, contraseña, datos personales
    \item Sistema valida datos y crea cuenta
    \item Sistema genera código QR único
\end{itemize}

\textbf{Precondición}: Miembro no registrado previamente

\textbf{Postcondición}: Cuenta creada, código QR generado
\end{definicion}

\subsubsection{CU-002: Realizar Check-in}

\begin{definicion}
\textbf{Descripción}: El miembro registra su entrada al gimnasio mediante QR.

\textbf{Flujo Principal}:
\begin{itemize}
    \item Actor presenta código QR en entrada
    \item Sistema escanea y valida QR
    \item Sistema registra fecha/hora de acceso
    \item Sistema otorga acceso a gimnasio
    \item Sistema actualiza leaderboard
\end{itemize}

\textbf{Alternativo}: Si QR inválido, denegar acceso y alertar

\textbf{Postcondición}: Asistencia registrada
\end{definicion}

\subsection{Actor: Entrenador}

\subsubsection{CU-010: Consultar Datos de Miembro}

\begin{definicion}
\textbf{Descripción}: El entrenador accede al perfil y historial de un miembro.

\textbf{Flujo Principal}:
\begin{itemize}
    \item Actor selecciona miembro de lista
    \item Sistema muestra perfil completo
    \item Sistema presenta historial de ejercicios
    \item Sistema muestra progresión (pesos, repeticiones)
\end{itemize}

\textbf{Precondición}: Entrenador autenticado

\textbf{Postcondición}: Datos mostrados
\end{definicion}

\subsubsection{CU-011: Recibir Recomendaciones de IA}

\begin{definicion}
\textbf{Descripción}: El entrenador recibe sugerencias de ejercicios del AI Copilot.

\textbf{Flujo Principal}:
\begin{itemize}
    \item Actor selecciona miembro
    \item Sistema analiza perfil y objetivos
    \item AI Service genera recomendaciones
    \item Sistema presenta opciones de ejercicios
\end{itemize}

\textbf{Postcondición}: Recomendaciones mostradas
\end{definicion}

\subsection{Actor: Gerente}

\subsubsection{CU-020: Visualizar Dashboard de Ocupación}

\begin{definicion}
\textbf{Descripción}: El gerente accede a dashboard en tiempo real de ocupación.

\textbf{Flujo Principal}:
\begin{itemize}
    \item Actor autenticado ingresa a dashboard
    \item Sistema obtiene datos de occupación actual
    \item Sistema muestra gráficos en tiempo real
    \item Sistema presenta alertas de anomalías
\end{itemize}

\textbf{Precondición}: Gerente autenticado

\textbf{Postcondición}: Dashboard actualizado
\end{definicion}

\subsubsection{CU-021: Generar Reportes}

\begin{definicion}
\textbf{Descripción}: El gerente genera reportes de operación.

\textbf{Flujo Principal}:
\begin{itemize}
    \item Actor selecciona tipo de reporte
    \item Sistema solicita parámetros (fechas, filtros)
    \item Sistema calcula datos
    \item Sistema exporta en PDF/Excel
\end{itemize}

\textbf{Postcondición}: Reporte generado y descargado
\end{definicion}

\subsection{Actor: Sistema}

\subsubsection{CU-030: Ejecutar Predicciones de Churn}

\begin{definicion}
\textbf{Descripción}: El sistema automáticamente detecta miembros en riesgo.

\textbf{Flujo Principal}:
\begin{itemize}
    \item Sistema analiza patrones de asistencia
    \item AI Service ejecuta modelo predictivo
    \item Sistema identifica usuarios en riesgo
    \item Sistema crea alertas para gerencia
\end{itemize}

\textbf{Frecuencia}: Diaria a las 00:00

\textbf{Postcondición}: Alertas generadas
\end{definicion}

% =========================================================
%                 CONCLUSIONES TÉCNICAS
% =========================================================
\section{Conclusiones Técnicas}

\subsection{Viabilidad Técnica}

\begin{itemize}
    \item \textbf{Stack}: Tecnologías maduras y probadas en producción
    \item \textbf{Escalabilidad}: Arquitectura microservicios permite crecimiento
    \item \textbf{Disponibilidad}: 99.9\% uptime alcanzable con cloud moderno
    \item \textbf{Seguridad}: Implementable cumpliendo estándares GDPR/CCPA
\end{itemize}

\subsection{Requisitos Satisfechos}

\begin{itemize}
    \item \textbf{Automatización}: 85\% de procesos automatizables
    \item \textbf{Múltiples plataformas}: Web, Mobile (iOS/Android), Desktop
    \item \textbf{Integraciones}: APIs de terceros para pagos, biometría, wearables
    \item \textbf{Análisis en tiempo real}: Analytics con IA implementable
\end{itemize}

\subsection{Complejidad del Proyecto}

\begin{table}[H]
\centering
\caption{Complejidad Técnica por Módulo}
\begin{tabular}{|l|l|r|}
\hline
\textbf{Módulo} & \textbf{Complejidad} & \textbf{Esfuerzo (meses)} \\ \hline
Member Service & Media & 1.5 \\ \hline
Attendance Service & Media & 1 \\ \hline
Payment Service & Alta & 2 \\ \hline
Analytics Service & Alta & 2.5 \\ \hline
AI Service & Alta & 3 \\ \hline
Notification Service & Media & 1 \\ \hline
Marketplace Service & Alta & 2 \\ \hline
Frontend Web & Alta & 2.5 \\ \hline
Frontend Mobile & Alta & 3 \\ \hline
\hline
\textbf{TOTAL} & & \textbf{18 meses} \\ \hline
\end{tabular}
\end{table}

\subsection{Riesgos Técnicos Identificados}

\begin{alerta}
\textbf{Riesgo Alto}:
\begin{itemize}
    \item Integración con múltiples wearables (variabilidad de APIs)
    \item Modelos de IA requieren data histórica suficiente
    \item Sincronización multi-sucursal en tiempo real
\end{itemize}

\textbf{Mitigación}:
\begin{itemize}
    \item Estándares de integración definidos (OpenAPI)
    \item MVP sin IA completa, agregar después
    \item Caché distribuida (Redis) para sincronización
\end{itemize}
\end{alerta}

\subsection{Próximas Fases}

\begin{enumerate}
    \item \textbf{Fase 3}: Especificar todos los RF y CU detallados
    \item \textbf{Fase 5}: Diseñar esquema de BD normalizado
    \item \textbf{Fase 6}: Crear wireframes y prototipo interactivo
    \item \textbf{Fase 7}: Desarrollo en sprints de 2 semanas
    \item \textbf{Testing}: QA y testing de penetración
    \item \textbf{Deployment}: Deployment a producción con CI/CD
\end{enumerate}

% =========================================================
%                   REFERENCIAS
% =========================================================
\newpage
\section*{Referencias}
\addcontentsline{toc}{section}{Referencias}

\begin{enumerate}
    \item DATOS\_PROYECTO.json (2026) — Control técnico del proyecto GIMNASIO

    \item CONSOLIDACION\_1\_2.md (2026) — Análisis integrado Fases 1+2

    \item FASE\_4\_PLAN\_NEGOCIO.md (2026) — Plan de negocio detallado

    \item Ries, E. (2011). \textit{The Lean Startup: How Today's Entrepreneurs Use Continuous Innovation to Create Radically Successful Businesses}. Crown Publishing.

    \item Blank, S., \& Dorf, B. (2012). \textit{The Startup Owner's Manual: The Step-By-Step Guide for Building a Great Company}. K\&S Ranch.

    \item Goodwin, T. (2015). \textit{The Platform Revolution: How Networked Markets Are Transforming the Economy and How to Make Them Work for You}. W.W. Norton \& Company.

    \item McKinsey \& Company (2024). \textit{The Future of Fitness: Digital Transformation in Health \& Wellness}.

    \item Gartner (2025). \textit{Market Guide for Fitness Management Software}.

    \item OpenAI Documentation (2025). \textit{GPT-4 API for Custom Applications}.
\end{enumerate}

% =========================================================
%                ANEXOS TÉCNICOS
% =========================================================
\newpage
\appendix

\section{Anexo A: Modelo de Datos Conceptual}

\subsection{Entidades Principales}

\begin{table}[H]
\centering
\caption{Entidades Principales del Sistema}
\begin{tabular}{|l|l|}
\hline
\textbf{Entidad} & \textbf{Descripción} \\ \hline
User & Usuario del sistema (miembro, entrenador, gerente) \\ \hline
Member & Información específica de miembros \\ \hline
Membership & Tipo y estado de membresía \\ \hline
Attendance & Registro de asistencias \\ \hline
Exercise & Ejercicios disponibles en el gimnasio \\ \hline
Workout & Planes de entrenamiento \\ \hline
Achievement & Logros y badges \\ \hline
Transaction & Transacciones de marketplace \\ \hline
Notification & Notificaciones enviadas \\ \hline
\end{tabular}
\end{table}

\subsection{Relaciones Principales}

\begin{itemize}
    \item User (1) → (N) Membership
    \item User (1) → (N) Attendance
    \item User (1) → (N) Achievement
    \item Member (1) → (N) Workout
    \item Workout (1) → (N) Exercise
    \item Membership (1) → (N) Transaction
\end{itemize}

\section{Anexo B: APIs del Sistema}

\subsection{REST API Endpoints Principales}

\begin{lstlisting}[caption={Endpoints REST del Member Service}]
POST   /api/v1/members           # Crear nuevo miembro
GET    /api/v1/members/{id}      # Obtener datos miembro
PUT    /api/v1/members/{id}      # Actualizar miembro
DELETE /api/v1/members/{id}      # Eliminar miembro
GET    /api/v1/members/{id}/history  # Historial de miembro

POST   /api/v1/membership        # Crear membresía
GET    /api/v1/membership/{id}   # Obtener membresía
PUT    /api/v1/membership/{id}   # Renovar membresía

POST   /api/v1/attendance        # Registrar asistencia
GET    /api/v1/attendance/{memberId} # Historial de asistencias
\end{lstlisting}

\section{Anexo C: Flujos de Datos Críticos}

\subsection{Flujo de Check-in}

\begin{lstlisting}[caption={Flujo de Check-in de Miembro}]
1. Cliente escanea QR
2. Mobile app valida QR
3. App envía POST /api/v1/attendance
4. Attendance Service valida miembro
5. Attendance Service registra timestamp
6. Analytics Service actualiza ocupación
7. Notification Service envía confirmación
8. Frontend muestra "Welcome!"
\end{lstlisting}

\subsection{Flujo de Predicción de Churn}

\begin{lstlisting}[caption={Flujo de Detección de Churn}]
1. Analytics Service ejecuta Churn AI diariamente
2. Obtiene datos de asistencia últimos 60 días
3. AI Service predice probabilidad de churn
4. Si probabilidad > 70%, crear alerta
5. Notification Service notifica a gerente
6. Notification Service sugiere intervención
7. Gerente recibe push notification y email
\end{lstlisting}

\section{Anexo D: Estándares de Codificación}

\subsection{Frontend}

\begin{itemize}
    \item Framework: React 18+
    \item Lenguaje: TypeScript
    \item Estado: Redux Toolkit
    \item Testing: Jest + React Testing Library
    \item Linting: ESLint + Prettier
\end{itemize}

\subsection{Backend}

\begin{itemize}
    \item Runtime: Node.js 18+
    \item Framework: Express.js
    \item Lenguaje: TypeScript
    \item Testing: Jest
    \item Linting: ESLint
\end{itemize}

\subsection{Base de Datos}

\begin{itemize}
    \item DBMS: PostgreSQL 14+
    \item ORM: Sequelize o TypeORM
    \item Migrations: Knex.js
    \item Caché: Redis
\end{itemize}

\section{Anexo E: Matriz de Rastreabilidad}

\begin{table}[H]
\centering
\caption{Rastreabilidad de Requisitos a Módulos}
\begin{tabular}{|l|l|l|}
\hline
\textbf{Requisito} & \textbf{Módulo} & \textbf{CU Relacionado} \\ \hline
RF-001 a RF-005 & Member Service & CU-001 \\ \hline
RF-010 a RF-014 & Attendance Service & CU-002 \\ \hline
RF-020 a RF-024 & Analytics Service & CU-020, CU-021 \\ \hline
RF-030 a RF-034 & Notification Service & - \\ \hline
RF-040 a RF-044 & Gamification Service & - \\ \hline
RF-050 a RF-054 & Marketplace Service & - \\ \hline
\end{tabular}
\end{table}

\end{document}
